\documentclass[11pt]{article}
\usepackage[margin=1.1in]{geometry}
\usepackage{amsmath,amssymb,amsthm}
\usepackage{mathtools}
\usepackage{enumitem}
\usepackage[colorlinks=true,linkcolor=blue,citecolor=blue]{hyperref}

\newtheorem{theorem}{Theorem}[section]
\newtheorem{lemma}[theorem]{Lemma}
\newtheorem{proposition}[theorem]{Proposition}
\newtheorem{corollary}[theorem]{Corollary}
\theoremstyle{definition}
\newtheorem{definition}[theorem]{Definition}
\newtheorem{algorithm}[theorem]{Algorithm}
\theoremstyle{remark}
\newtheorem{remark}[theorem]{Remark}

\newcommand{\R}{\mathbb{R}}
\newcommand{\T}{\boldsymbol{T}}
\newcommand{\rank}{\operatorname{rank}}
\newcommand{\tr}{\operatorname{tr}}
\newcommand{\lammax}{\lambda_{\max}}
\newcommand{\lammin}{\lambda_{\min}}
\newcommand{\HS}{\mathrm{HS}}
\newcommand{\cond}{\operatorname{cond}}
\newcommand{\diag}{\operatorname{diag}}
\newcommand{\Mnr}{\mathcal{M}_{\boldsymbol{n},\boldsymbol{r}}}
\newcommand{\Msh}{\mathcal{M}^{\mathrm{sh}}_{\boldsymbol{n},\boldsymbol{r}}}

\title{Shared Tucker Factors for Tucker Tensor Trains:\\
Algorithms, Bounds, and Proofs}
\author{Companion notes for the T3Toolbox sharing implementation}
\date{\today}

\begin{document}
\maketitle

\begin{abstract}
\noindent
These notes summarize the mathematics behind adding \emph{shared Tucker
factors} to the Tucker tensor train (T3) format: the shared rank theorem, the
grouped truncation algorithm and its quasi-optimality bound, initialization
from unshared representations, the geometry of the fixed-rank shared manifold
(tied tangent spaces, the Gram-sum projection, retraction, dimension), the
edge-conditioning statistics used for rank continuation, and the single-pass
property of group-aware useless-rank removal. The shared format is the SF-ETT
decomposition of Molozhavenko and Rakhuba~\cite{MR2026}, generalized from one
trailing shared block to an arbitrary partition of the modes into sharing
groups.
\end{abstract}

\noindent\emph{Revision (v2, 2026-08-19): updated after implementation review
--- Proposition~\ref{prop:tilt} restated basis-free (no complement basis),
the corewise geometry added (Proposition~\ref{prop:corewise}), the
numerical-form and zero-padded-restart remarks added
(Remarks~\ref{rem:svdsolve} and~\ref{rem:restart}), the dimension formula's
limits corrected (Proposition~\ref{prop:dim} and its erratum remark), the
tied doubled-rank embedding made explicit (Remark~\ref{rem:embedding}), and
the tail-energy definition of $\rank_\varepsilon$ made explicit
(Theorem~\ref{thm:rankbound}).}

\section{Setup and notation}\label{sec:setup}

A \emph{Tucker tensor train} (T3; equivalently, extended tensor train)
represents a tensor $\T \in \R^{N_1 \times \cdots \times N_d}$ as
\[
\T = \bigl[\bigl[\, G ;\; U_1, \ldots, U_d \,\bigr]\bigr],
\qquad
G = \mathrm{TT}(G_1, \ldots, G_d),
\]
where $U_i \in \R^{N_i \times n_i}$ are the Tucker factors,
$G_i \in \R^{r_{i-1} \times n_i \times r_i}$ are the TT cores of the central
Tucker core $G \in \R^{n_1 \times \cdots \times n_d}$, and
$r_0 = r_d = 1$. We call $\boldsymbol{n} = (n_1,\ldots,n_d)$ the Tucker ranks
and $\boldsymbol{r} = (r_1,\ldots,r_{d-1})$ the TT ranks.

$\T_{(i)} \in \R^{N_i \times \prod_{j\neq i} N_j}$ denotes the $i$-th
matricization of $\T$ and $\T^{\langle i \rangle}$ its $i$-th matrix
unfolding. $W_i \in \R^{n_i \times r_{i-1} r_i}$ denotes the
\emph{up-matricization} of core $G_i$ (Tucker leg to rows). All statements
below assume the standard \emph{orthogonalized representations}: the factors
$U_i$ have orthonormal columns, and, when a statement refers to mode $i$, the
TT chain is in \emph{$i$-orthogonal form} (cores left of $i$ left-orthogonal,
cores right of $i$ right-orthogonal), reachable by QR sweeps without changing
$\T$.

\begin{definition}[Mode environment / Gram matrix]\label{def:gram}
$\Gamma_i \coloneqq G_{(i)} G_{(i)}^\top \in \R^{n_i \times n_i}$, the Gram of
the $i$-th matricization of the Tucker core. $\Gamma_i$ is
representation-independent; in the $i$-orthogonal representation it
telescopes to
\[
\Gamma_i = W_i W_i^\top .
\]
We write $S_i$ for any factor with $\Gamma_i = S_i S_i^\top$ (e.g.\ the
R-factor absorbed during mode-$i$ down-orthogonalization; in frame coordinates
of the fixed-rank manifold, $S_i$ is pre-absorbed so that the coordinate Gram
is the identity).
\end{definition}

\begin{definition}[Sharing partition; shared T3]\label{def:sharing}
A \emph{sharing partition} assigns each mode $i \in \{1,\ldots,d\}$ to a group
$g$; modes in a group $g = (i_1, \ldots, i_{k_g})$ must have equal mode sizes,
denoted $N_g$. A T3 is \emph{shared} (an SF-T3) if $U_{i} = U_{g}$ and
$n_i = n_g$ for all $i \in g$, for every group. Singleton groups impose no
constraint. $\Msh$ denotes the set of shared T3s with fixed ranks
(Section~\ref{sec:geometry}).
\end{definition}

Throughout, ``the concatenation'' for a group $g$ means the column-wise block
matrix $[\,\T_{(i_1)} \,|\, \cdots \,|\, \T_{(i_{k_g})}\,]$.

\section{Shared ranks}\label{sec:ranks}

\begin{lemma}[Stacked skeleton factorization]\label{lem:skeleton}
Let $\T$ be a shared T3 with orthonormal factors, and let each $W_i$ be taken
in the $i$-orthogonal representation. Then for each mode $i$ of group $g$,
\begin{equation}\label{eq:skeleton}
\T_{(i)} = U_g\, W_i\, M_i,
\end{equation}
where $M_i$ has orthonormal rows, and stacking the group's matricizations,
\begin{equation}\label{eq:stacked}
\bigl[\,\T_{(i_1)} \,|\, \cdots \,|\, \T_{(i_k)}\,\bigr]
= U_g \,\bigl[\, W_{i_1} \,|\, \cdots \,|\, W_{i_k} \,\bigr]\,
\operatorname{blockdiag}\!\bigl(M_{i_1}, \ldots, M_{i_k}\bigr).
\end{equation}
Consequently the nonzero singular values of the concatenation equal those of
the small matrix $[\,W_{i_1} | \cdots | W_{i_k}\,]$, and
\begin{equation}\label{eq:specid}
s_{g,j}^2 \;=\;
\lambda_j\Bigl(\textstyle\sum_{i \in g} \Gamma_i\Bigr)
\;=\;
\lambda_j\Bigl(\textstyle\sum_{i \in g} S_i S_i^\top\Bigr),
\qquad j = 1, \ldots, n_g,
\end{equation}
where $s_{g,1} \geq s_{g,2} \geq \cdots$ are the singular values of the
concatenation.
\end{lemma}

\begin{proof}
Equation~\eqref{eq:skeleton} is the standard matricization identity for a T3
in $i$-orthogonal form: the environment of mode $i$'s Tucker leg contracts to
a matrix $M_i$ with orthonormal rows (a Kronecker/interface product of
orthonormal-column factors and left/right-orthogonal partial chains).
Each block of the concatenation is factored in its own representation of the
same tensor $\T$; since each factorization holds individually and the left
factor $U_g$ is common by sharedness, \eqref{eq:stacked} follows by placing
the $M_i$ in block-diagonal position. The outer factors in \eqref{eq:stacked}
are, respectively, orthonormal-column and orthonormal-row, so singular values
pass to the middle factor. Finally,
$[\,W_{i_1}|\cdots|W_{i_k}\,][\,W_{i_1}|\cdots|W_{i_k}\,]^\top
= \sum_{i\in g} W_i W_i^\top = \sum_{i\in g} \Gamma_i$, and $\Gamma_i$ is
representation-independent, so mixing per-mode representations is consistent.
\end{proof}

\begin{theorem}[Shared minimal rank; cf.~{\cite[Thm.~1]{MR2026}},
{\cite{Pesh2024}}]\label{thm:sharedrank}
The minimal Tucker rank of group $g$ over all shared T3 representations of
$\T$ is
\[
n_g^{\min} \;=\; \rank\bigl[\,\T_{(i_1)} \,|\, \cdots \,|\, \T_{(i_k)}\,\bigr].
\]
TT ranks are unaffected by sharing:\ $r_i^{\min} = \rank \T^{\langle i\rangle}$.
\end{theorem}

\begin{proof}[Proof sketch]
$(\geq)$: in any shared representation, every $\T_{(i)}$, $i \in g$, has
column space contained in $\operatorname{range}(U_g)$, hence so does the
concatenation, giving $\rank[\cdots] \leq n_g$.
$(\leq)$: take $U_g$ to be an orthonormal basis of the concatenation's column
space; projecting each mode ($U_g U_g^\top \T_{(i)} = \T_{(i)}$) yields an
exact shared representation with $n_g = \rank[\cdots]$; TT-decompose the
resulting core at the unfolding ranks. See \cite{MR2026,Pesh2024} for
details.
\end{proof}

\section{Grouped truncation (shared rounding)}\label{sec:rounding}

\begin{algorithm}[Grouped T3-SVD / SF-rounding; cf.~{\cite[Alg.~1]{MR2026}}]
\label{alg:rounding}
Input: shared T3 with orthonormal factors; target ranks. Output: shared T3
with target ranks.
\begin{enumerate}[label=\emph{(\arabic*)},itemsep=1pt,topsep=2pt]
\item TT-round the cores to the target TT ranks (standard sweep; factors
  untouched). In an interlaced implementation, steps (2)--(4) are woven into
  the same sweep; the interlacing order changes truncated results but not the
  bounds below.
\item For each mode $i$, obtain $W_i$ in the $i$-orthogonal representation
  (one orthogonalization sweep).
\item Singleton modes: truncated SVD of $W_i$; rotate factor and core as
  usual.
\item Each group $g$: one truncated SVD of the concatenation
  $[\,W_{i_1} | \cdots | W_{i_k}\,] \in \R^{n_g \times \sum_{i\in g}
  r_{i-1} r_i}$, yielding leading left singular vectors $Y \in \R^{n_g \times
  n_g'}$. Update $U_g \leftarrow U_g Y$ \emph{once} (assigned to every mode of
  $g$) and apply $Y^\top$ to the up-leg of every core in $g$.
\end{enumerate}
By Lemma~\ref{lem:skeleton}, step (4) computes the optimal rank-$n_g'$
truncation of the concatenation while touching the large dimension $N_g$ only
in the single product $U_g Y$, at cost $O(N_g n_g^2)$ per group; the SVD
itself acts on an $n_g \times \sum_i r_{i-1} r_i$ matrix.
\end{algorithm}

\begin{theorem}[Quasi-optimality; {\cite[Thm.~2]{MR2026}}]\label{thm:quasi}
Let $P_{\mathrm{sf}}(\T)$ denote the output of
Algorithm~\ref{alg:rounding} (equivalently, of its dense analogue applied to
an arbitrary tensor) at ranks $(\boldsymbol{n}, \boldsymbol{r})$. Then
\[
\bigl\| \T - P_{\mathrm{sf}}(\T) \bigr\|_{\HS}
\;\leq\;
C(d)\,
\inf_{\substack{\boldsymbol{Y} \text{ shared},\\
\mathrm{ranks}(\boldsymbol{Y}) \preceq (\boldsymbol{n},\boldsymbol{r})}}
\bigl\| \T - \boldsymbol{Y} \bigr\|_{\HS},
\qquad
C(d) = \sqrt{d} + \sqrt{d}\sqrt{d-1} + \sqrt{d-1}.
\]
\end{theorem}

\begin{proof}[Proof sketch]
Split $P_{\mathrm{sf}} = P_{\mathrm{SF\text{-}Tucker}} \circ
P_{\mathrm{TT}}$ (either order; \cite[Rem.~1]{MR2026}). Each stage is
quasi-optimal in its own class ($\sqrt{d}$ for the HOSVD-type shared factor
truncation, using that the shared-mode error is controlled by the
concatenation tail; $\sqrt{d-1}$ for TT-SVD), and each stage's best-in-class
error is bounded by the best shared error via the triangle inequality; the
constants compose to $C(d)$.
\end{proof}

\begin{remark}[Truncation error bound]
When the input is exactly shared and step (4) truncates group $g$ from
$n_g$ to $n_g'$ with retained basis $Y$, the squared error contributed by
that step is \emph{at most} $\sum_{j > n_g'} s_{g,j}^2$: by
Lemma~\ref{lem:skeleton} the discarded tail of the concatenation equals the
\emph{sum over the group's modes} of the single-mode projection errors,
$\sum_{i \in g} \|(I - YY^\top)\T_{(i)}\|_F^2$, which bounds the error of the
simultaneous projection $\T \times_{i \in g} YY^\top$ by the standard
HOSVD subadditivity argument. Equality holds for singleton groups ($k=1$);
for $k > 1$ the simultaneous projection error is strictly smaller in general
(verified numerically: e.g.\ $4.19\cdot10^2$ vs the tail bound
$4.80\cdot10^2$ on a random tied example).
\end{remark}

\begin{theorem}[Rank upper bound for the interlaced sweep]\label{thm:rankbound}
In tolerance-based truncation, every rank selected by
Algorithm~\ref{alg:rounding} satisfies
\[
n_g' \;\leq\; \rank_\varepsilon
\bigl[\,\T_{(i_1)} \,|\, \cdots \,|\, \T_{(i_k)}\,\bigr]
\quad\text{(of the \emph{original} input)},
\qquad
r_i' \;\leq\; \rank_\varepsilon \T^{\langle i \rangle},
\]
where $\rank_\varepsilon$ is numerical rank at the working tolerance,
counted by tail Frobenius energy (the criterion the implemented truncated SVD
uses): $\rank_\varepsilon(A) = \min\{\, r :
\|A - A_r\|_F \le \max(\varepsilon \|A\|_F,\ \mathrm{atol}) \,\}$ with $A_r$
the best rank-$r$ approximation.
\end{theorem}

\begin{proof}[Proof sketch]
Edge monotonicity: any truncation performed at another edge multiplies the
group's concatenation on the right by a block-orthogonal-projection-type
factor (and on the left by nothing, since factors are orthonormal), which
cannot increase its rank; likewise for unfoldings. Hence the matrix whose SVD
is truncated at each step has rank at most that of the corresponding
matricization/unfolding of the original tensor.
\end{proof}

\section{Initialization: unshared $\to$ shared}\label{sec:init}

\begin{algorithm}[Shared projection of an unshared T3]\label{alg:init}
Input: T3 with orthonormal but distinct factors $U_i$, $i \in g$.
For each group $g$:
\begin{enumerate}[label=\emph{(\arabic*)},itemsep=1pt,topsep=2pt]
\item Stack factors and compute one tall QR:
  $[\,U_{i_1} | \cdots | U_{i_k}\,] = Q R$, with
  $Q \in \R^{N_g \times m}$, $m \leq \sum_{i\in g} n_i$.
\item Express each block $B_i \coloneqq U_i W_i$ (never formed at size
  $N_g$) in the $Q$ basis: $Q^\top B_i = R_{[i]} W_i$, where $R_{[i]}$ is the
  corresponding column block of $R$. SVD the small stacked matrix
  $[\,R_{[i_1]} W_{i_1} | \cdots | R_{[i_k]} W_{i_k}\,] \in
  \R^{m \times \sum_i r_{i-1} r_i}$; retain leading left singular vectors
  $Y$.
\item Set $U_g \coloneqq Q Y$; project each mode by applying
  $(U_g^\top U_i)$ to the up-leg of core $i$ and replacing the factor by
  $U_g$. Follow with Algorithm~\ref{alg:rounding}.
\end{enumerate}
\end{algorithm}

\begin{lemma}\label{lem:init}
The column space of $[\,B_{i_1} | \cdots | B_{i_k}\,]$ lies in
$\operatorname{range}(Q)$, so step (2) computes the exact dominant left
singular subspace of the concatenation $[\,\T_{(i_1)}|\cdots|\T_{(i_k)}\,]$
while touching the large dimension only in the QR of step (1) and the product
$QY$. The resulting projection is quasi-optimal with respect to the best
shared approximation, with the constant of Theorem~\ref{thm:quasi} (same
composition argument, with the shared-HOSVD stage now acting on unshared
input).
\end{lemma}

\begin{remark}[Implementation form: exact rewrite, defer all selection to the
rounding]
T3Toolbox implements only step (1)'s stacked-factor SVD, as an \emph{exact}
common-span rewrite: the factor coefficients are the SVD's own row blocks
($B_i = W_{[i]} \operatorname{diag}(s)\, V^\top$, valid for arbitrary --- not
necessarily orthonormal --- input factors), the shared factor becomes
$V^\top$ (group rank $m = \min(\sum_i n_i, N_g)$), and ALL rank/tolerance
selection is deferred to the grouped rounding
(Algorithm~\ref{alg:rounding}). This is equivalent to steps (2)--(3): the
optimal shared basis lies in the span of the group's factors, and the group
spectrum is representation-independent, so the rounding on the rewritten
representation selects the same subspace and reports the same $s_g$; the
explicit small SVD of steps (2)--(3) is an optional efficiency device
(pre-shrinking the group rank before the TT sweep), at the cost of a second
truncation stage.
\end{remark}

\section{Fixed-rank geometry with tied factors}\label{sec:geometry}

Let $\Mnr$ denote the manifold of fixed-rank T3s and
$T_{\T}\Mnr$ its tangent space at a point $\T$ with orthonormal
representation. A tangent vector is a sum of $d$ TT-variation terms and $d$
Tucker-variation terms; in \emph{gauged} form the Tucker variation at mode
$i$ satisfies $U_i^\top \dot V_i = 0$, and gauge conditions make all cross
terms in the Hilbert--Schmidt inner product vanish, so
\begin{equation}\label{eq:hsnorm}
\langle \xi, \eta \rangle_{\HS}
= \sum_{i=1}^{d} \tr\bigl( \dot V_i^\top \dot W_i\, \Gamma_i \bigr)
+ \text{(TT-variation terms)} .
\end{equation}
In \emph{frame coordinates} the factor $S_i$ (with $\Gamma_i = S_i S_i^\top$)
is pre-absorbed per mode, so the coordinates $C_i$ of Tucker variations are
orthonormal: corewise Euclidean inner products equal Hilbert--Schmidt inner
products.

\begin{definition}[Shared manifold]
$\Msh \subset \Mnr$ is the set of shared T3-representable tensors with fixed
shared ranks: one Tucker rank $n_g$ and one factor per group. It is the
transversal intersection of the fixed-rank T3 manifold with the shared-Tucker
manifold and is smooth \cite[Thm.~5]{MR2026}.
\end{definition}

\begin{proposition}[Tied tangent space is a tilted subspace]\label{prop:tilt}
Fix a shared point. A tangent vector is tangent to $\Msh$ iff its TT
variations are unrestricted and, per group, its Tucker variations arise from
a single \emph{gauged ambient direction} $\dot U \in \R^{N_g \times n_g}$
with $U_g^\top \dot U = 0$: in the ($S$-absorbed) frame coordinates,
\begin{equation}\label{eq:tied}
C_i = \dot U\, S_i \quad \text{for a common gauged } \dot U, \qquad i \in g .
\end{equation}
This is a linear subspace of the full gauged tangent space, tilted
mode-by-mode by the (generally distinct, generally rectangular) $S_i$.
\end{proposition}

\begin{remark}[No complement basis; rectangular $S_i$]
An earlier draft stated \eqref{eq:tied} through an orthonormal complement
$O_g$ of $U_g$ ($\dot U = O_g B$, coordinates $C_i = B S_i$), suggesting a
complement must be constructed once per group and tied across it. The
subspace depends only on $U_g$ (tied by hypothesis) and the $S_i$; no
complement basis is ever materialized (nor is one materialized anywhere in
T3Toolbox), and the per-mode rotational freedom in the factorization
$\Gamma_i = S_i S_i^\top$ (replacing $S_i$ by $S_i Q_i$ together with the
correspondingly rotated orthonormal core) cancels identically in
\eqref{eq:tied} and in the projection \eqref{eq:gramsolve}. Note also that in
implementation $S_i \in \R^{n_g \times n_{D,i}}$ with
$n_{D,i} = \min(n_g,\ r_{i-1} r_i)$ --- rectangular in general, so an
individual Gram $\Gamma_i$ may be singular at a perfectly smooth shared
point; only the group sum is invertible there
(Proposition~\ref{prop:kappa}(iii)).
\end{remark}

\begin{proposition}[Projection onto the tied subspace]\label{prop:proj}
The orthogonal projection (in the frame-coordinate Euclidean metric, which
equals Hilbert--Schmidt) of arbitrary gauged coordinates
$(C_{i})_{i \in g}$ onto the subspace~\eqref{eq:tied} is given by
\begin{equation}\label{eq:gramsolve}
\dot U \;=\; \Bigl(\sum_{i \in g} C_i S_i^\top\Bigr)
        \Bigl(\sum_{i \in g} S_i S_i^\top\Bigr)^{-1}
   \;=\; \Bigl(\sum_{i \in g} C_i S_i^\top\Bigr)
        \Bigl(\sum_{i \in g} \Gamma_i\Bigr)^{-1},
\qquad
C_i \leftarrow \dot U S_i .
\end{equation}
TT variations and singleton modes are unaffected, and $\dot U$ is
automatically gauged when the $C_i$ are. Formula~\eqref{eq:gramsolve}
is the projection formula of \cite[\S 5.3]{MR2026} rederived as least squares
onto a tilted diagonal; $\sum_i C_i S_i^\top$ equals the sum of raw per-mode
cotangents $(\nabla f)_{(i)} G_{(i)}^\top$.
\end{proposition}

\begin{proof}
Minimize $\sum_{i\in g} \| C_i - \dot U S_i \|_F^2$ over $\dot U$; setting the
gradient to zero gives the normal equations
$\sum_i C_i S_i^\top = \dot U \sum_i S_i S_i^\top$, invertible whenever the
point has full shared rank ($s_{g,n_g} > 0$ by \eqref{eq:specid}). Gauge:
$U_g^\top \sum_i C_i S_i^\top = \sum_i (U_g^\top C_i) S_i^\top = 0$, and the
Gram sum acts on the other side, so $U_g^\top \dot U = 0$.
\end{proof}

\begin{remark}[Numerical form: solve the tilted least squares by SVD, never
by normal equations]\label{rem:svdsolve}
Implementations should not form the Gram sum. Stacking
\[
M_g = \begin{bmatrix} S_{i_1}^\top \\ \vdots \\ S_{i_k}^\top \end{bmatrix}
\in \R^{(\sum_{i\in g} n_{D,i}) \times n_g},
\qquad
C^{\mathrm{stack}} = \begin{bmatrix} C_{i_1}^\top \\ \vdots \\
C_{i_k}^\top \end{bmatrix},
\]
the projection is the least-squares solution
$\dot U^\top = M_g^{+}\, C^{\mathrm{stack}}$ (then
$C_i^\top \leftarrow S_i^\top \dot U^\top$, the corresponding row block of
$M_g \dot U^\top$), computed from one thin SVD of $M_g$ with a relative
singular-value clip. Three payoffs. (i) The solve has the intrinsic
least-squares sensitivity ($\kappa_g$ on the solution component); the
normal-equations/Cholesky route is $\kappa_g^2$ unconditionally.
(ii) $\sigma(M_g) = (s_{g,j})_j$ \emph{is} the group spectrum
\eqref{eq:specid}, at full (non-squared) accuracy: in float32 the Gram route
loses trailing levels below $\sqrt{\varepsilon}\, s_{g,1}$ entirely
(measured at $s_{g,\min}/s_{g,1} \approx 10^{-4}$: relative error
$3\cdot10^{-1}$ via the Gram eigendecomposition vs $8\cdot10^{-5}$ via the
SVD of $M_g$). (iii) The clipped pseudoinverse is well-defined
(minimum-norm) at rank-deficient points, which rank continuation visits by
construction (Remark~\ref{rem:restart}).
\end{remark}

\begin{remark}[Metric unchanged]
Because \eqref{eq:tied} is a linear subspace and the frame metric is
Euclidean on gauged coordinates, the restriction of the metric to tied
tangents requires no correction: inner products and norms of tied tangents
are computed exactly as in the unshared case. The Gram-sum enters only
through the projection~\eqref{eq:gramsolve}. In particular, plain averaging
$C_i \mapsto \frac1k \sum_j C_j$ projects onto the wrong subspace
$\{C_{i_1} = \cdots = C_{i_k}\}$, which corresponds to no tied ambient
direction.
\end{remark}

\begin{proposition}[The corewise geometry ties by averaging]
\label{prop:corewise}
Consider instead the over-parametrized \emph{corewise} geometry: points are
the raw cores $((U_i)_i, (G_i)_i)$, the metric is Euclidean on the core
entries, and the retraction is addition of core perturbations. The shared
parameter set $\{\,U_i = U_j,\ i, j \in g,\ \forall g\,\}$ is a linear
subspace of the parameter space; its tangent space is
$\{\dot U_i \text{ equal within each group}\}$ (TT-core perturbations
unrestricted), and the orthogonal projection onto it in the corewise metric
is the per-group arithmetic mean
$\dot U_i \mapsto \tfrac{1}{k_g} \sum_{j \in g} \dot U_j$. The additive
retraction then preserves tying exactly.
\end{proposition}

\begin{remark}[One principle, two formulas]
Both geometries tie by \emph{orthogonal projection onto the geometry's tied
subspace in the geometry's own metric on its own coordinates}. The formulas
differ because the coordinates differ: the manifold geometry's gauged
coordinates carry $S_i$ pre-absorbed, so its tied subspace is tilted
\eqref{eq:tied} and the projection is the Gram-weighted solve
\eqref{eq:gramsolve}; the corewise coordinates are raw factor copies, so its
tied subspace is the diagonal and the projection is the mean. Using either
formula in the other geometry projects onto the wrong subspace; applied to
the same gauged coordinates, the arithmetic mean matches the Gram-weighted
solve only when the $S_i$ themselves coincide across the group (equal Grams
$\Gamma_i$ are not sufficient -- the $S_i$ are pinned to the per-mode
orthonormal cores and generally differ by more than a common rotation even
at exactly symmetric points).
\end{remark}

\begin{proposition}[Retraction]\label{prop:retract}
For a shared point $\T$ and tied tangent $\xi$, the attached point
$\T + \xi$ admits a shared T3 representation with doubled ranks (the doubled
factor is $[\,U_g \,|\, \dot U_g\,]$, identical across each group), and
\[
R(\T, \xi) \coloneqq P_{\mathrm{sf}}(\T + \xi)
\]
(Algorithm~\ref{alg:rounding} truncating back to the original ranks) is a
retraction on $\Msh$. This follows from quasi-optimality
(Theorem~\ref{thm:quasi}) via the argument of
\cite[Prop.~2.3]{KSV2014}; see \cite[Prop.~6]{MR2026}.
\end{proposition}

\begin{remark}[Tied doubled-rank embedding]\label{rem:embedding}
The attached point $\T + \xi$ should be \emph{constructed} tied, not re-tied
after the fact: per group, the doubled factor is $[\,U_g \,|\, \dot U\,]$
(one matrix, used at every mode of $g$; $2 n_g$ columns) and the paired core
block at mode $i$ is the \emph{center core} $W_i$ (equivalently, $S_i$
absorbed back into the $S$-stripped orthonormal core) --- cf.\ the SF-ETT
tangent representation \cite[\S 5.2]{MR2026}. The naive embedding by
per-mode factor blocks $[\,U_g \,|\, \dot U S_i\,]$ (the $S$-absorbed
coordinates) is \emph{not} tied: the blocks differ across the group in value
and, when the $n_{D,i}$ differ, even in shape, and no per-mode averaging
repairs this. $\dot U$ is recovered from tied coordinates by the
least-squares solve of Remark~\ref{rem:svdsolve} (exact on tied input; its
residual is a natural tied-tangent check).
\end{remark}

\begin{proposition}[Dimension; cf.~{\cite[Thm.~5]{MR2026}}]\label{prop:dim}
With groups $g$ of mode size $N_g$ and shared rank $n_g$ (singleton groups
reducing to the standard terms),
\[
\dim \Msh
= \underbrace{\sum_{i=1}^{d} r_{i-1}\, n_{g(i)}\, r_i
  \;-\; \sum_{i=1}^{d-1} r_i^2}_{\text{TT-core term (cores are never tied)}}
\;+\; \sum_{\text{groups } g} n_g \bigl( N_g - n_g \bigr),
\]
where $g(i)$ is the group of mode $i$ and the boundary bonds
$r_0 = r_d = 1$ carry no gauge: one Stiefel-quotient term per \emph{group}
instead of per mode.
\end{proposition}

\begin{remark}[Erratum and scope]
The printed \cite[Thm.~5]{MR2026} subtracts $\sum_{i=1}^{d} r_i^2$,
over-subtracting by exactly $1$ (the boundary bond $r_d = 1$ is not a gauge
degree of freedom); the matrix case $d = 2$ pins the limits above. Note also
that both \cite{MR2026} and \cite{Pesh2024} prove smoothness and the
dimension for a \emph{single trailing block} of shared modes; the
arbitrary-partition statement here follows by iterating the same
transversality argument group by group, and has been validated empirically by
matching the rank of a dense tied tangent basis to the formula --- including
a case where a group rank exceeds an individual mode's local ceiling
$r_{i-1} r_i$ (Proposition~\ref{prop:ceiling}), where the \emph{unshared}
minimal-rank reduction would clip the rank and the formula, applied to the
unshared-reduced ranks, under-counts. Implementations must therefore reduce
by the \emph{shared} minimal ranks before applying the formula.
\end{remark}

\section{Rank continuation: edge conditioning under sharing}\label{sec:cont}

The rank-continuation policy of \cite{T4S} grows ranks on the most
well-conditioned edges of the tensor network, keeping all edges comparably
conditioned, because edge conditioning governs the conditioning of the
fixed-rank subproblem. Sharing replaces the group's per-mode Tucker edges
with a single edge whose spectrum is the concatenation spectrum
$s_{g,\cdot}$ of \eqref{eq:specid}:
\[
\kappa_g \coloneqq \frac{s_{g,1}}{s_{g,n_g}},
\qquad
\kappa^{\mathrm{loc}}_i \coloneqq \sqrt{\cond(\Gamma_i)}
\quad (i \in g),
\]
competing in the same growth rule as TT-edge and singleton-Tucker condition
numbers: an edge grows only if it is well conditioned relative to the worst
edge, $\kappa < \kappa_{\max}/\tau$ (with $\tau > 1$); the group contributes
\emph{one} $\kappa_g$ to the pool, and one growth decision is applied
group-wide.

\begin{proposition}[$\kappa_g$ is the shared subproblem conditioning]
\label{prop:kappa}
At a shared point:
\begin{enumerate}[label=(\roman*),itemsep=1pt,topsep=2pt]
\item The Gram of the tied parameterization $B \mapsto (BS_i)_{i \in g}$ is
  $\sum_{i\in g}\Gamma_i$: for a tied Tucker tangent,
  $\|\xi\|_{\HS}^2 = \tr\bigl(B (\sum_i \Gamma_i) B^\top\bigr)$. Its
  eigenvalue spread is $\kappa_g^2$.
\item The projection solve~\eqref{eq:gramsolve} has condition number
  $\kappa_g^2$.
\item $s_{g,n_g} \to 0$ iff the point approaches the stratum of lower shared
  rank (Theorem~\ref{thm:sharedrank}), where the fixed-rank parametrization
  degenerates. Individual $\Gamma_i$ are never inverted for shared modes;
  per-mode near-rank-deficiency is not a degeneracy of $\Msh$.
\end{enumerate}
\end{proposition}

\begin{proof}
(i) is \eqref{eq:hsnorm} with $\dot V_i \leftrightarrow BS_i$ absorbed into
frame coordinates: $\sum_i \|BS_i\|_F^2 = \tr(B(\sum_i S_iS_i^\top)B^\top)$.
(ii) is immediate from \eqref{eq:gramsolve} and \eqref{eq:specid}. (iii) is
Theorem~\ref{thm:sharedrank} with \eqref{eq:specid}.
\end{proof}

\begin{proposition}[Conditioning of the sum]\label{prop:mediant}
$\displaystyle \kappa_g \;\leq\; \max_{i \in g} \kappa^{\mathrm{loc}}_i$,
and the inequality can be strict by an arbitrarily large factor.
\end{proposition}

\begin{proof}
All $\Gamma_i$ are positive semidefinite, so
$\lammax(\sum_i \Gamma_i) \leq \sum_i \lammax(\Gamma_i)$ and
$\lammin(\sum_i \Gamma_i) \geq \sum_i \lammin(\Gamma_i)$. Hence, by the
mediant inequality
$\frac{\sum_i a_i}{\sum_i b_i} \leq \max_i \frac{a_i}{b_i}$
(for $a_i \ge 0$, $b_i > 0$),
\[
\kappa_g^2
= \frac{\lammax(\sum_i \Gamma_i)}{\lammin(\sum_i \Gamma_i)}
\leq \frac{\sum_i \lammax(\Gamma_i)}{\sum_i \lammin(\Gamma_i)}
\leq \max_{i \in g} \frac{\lammax(\Gamma_i)}{\lammin(\Gamma_i)}
= \max_{i \in g} \cond(\Gamma_i).
\]
Strictness: $\Gamma_1 = \diag(1, \epsilon)$, $\Gamma_2 = \diag(\epsilon, 1)$
gives $\kappa^{\mathrm{loc}}_1 = \kappa^{\mathrm{loc}}_2 = \epsilon^{-1/2}$
but $\kappa_g = 1$.
\end{proof}

Proposition~\ref{prop:mediant} says the shared edge is never worse
conditioned than the worst per-mode Tucker edge of its group, and can be far
better: under tying, a direction is well-determined if \emph{some} mode of
the group informs it. Conversely, $\kappa_g$ remaining small while $n_g$
grows past each mode's individual effective rank signals a shared basis being
forced to cover a union of dissimilar column spaces (sharing a poor model for
the data).

\begin{corollary}[Scale cancellation; exact symmetric degeneration]
\label{cor:sym}
Each $i$-orthogonal representation preserves norms, so
$\sum_j s_{g,j}^2 = k_g \|\T\|_{\HS}^2$: raw magnitudes inflate by
$\sqrt{k_g}$ relative to a single mode but cancel in the ratio $\kappa_g$,
which compares against TT and singleton edges with no normalization. If the
group is exactly symmetric ($\Gamma_{i_1} = \cdots = \Gamma_{i_k} =
\Gamma$; e.g.\ symmetric derivative tensors with all input modes shared),
then $\sum_i \Gamma_i = k\,\Gamma$, so $s_{g,j} = \sqrt{k}\,\sigma_j$
elementwise and $\kappa_g$ equals the per-mode condition number exactly.
\end{corollary}

\begin{remark}[Zero-padded restarts: which tangent channels are alive]
\label{rem:restart}
Rank continuation warm-starts by zero-padding, which places the iterate
\emph{on} the lower-shared-rank stratum: the freshly added factor directions
carry no core mass, so the corresponding rows of every $S_i$ vanish, the new
levels of $s_g$ are exactly zero, and the tied Tucker channel
\eqref{eq:tied} admits no first-order motion in the new directions (the
clipped pseudoinverse of Remark~\ref{rem:svdsolve} returns the minimum-norm
solution --- zero there). The escape is carried by the \emph{untied
TT-variation channel}: TT variations are unrestricted and have first-order
content in the new up-slots, paired with the (deterministically completed,
tied) new factor columns, so the first Gauss--Newton step activates the new
directions through core mass; once $s_g$ acquires mass, the tied Tucker
channel reactivates, with a transiently large $\kappa_g$ that the clipped
solve handles (small mass $\Rightarrow$ large permitted rotation --- the
same behavior the unshared $S$-absorbed coordinates exhibit). Consequently,
full shared rank must \emph{not} be enforced as a precondition of the
projection, and the projection solve must be a clipped pseudoinverse. All
statements in this remark have been verified numerically on zero-padded
shared points.
\end{remark}

\section{Group-aware useless-rank removal: single pass suffices}
\label{sec:removal}

Rank proposals are reduced to nondegenerate values by the standard
three-phase sweep (Tucker vs.\ mode size; TT right-unfoldings right-to-left;
interleaved Tucker/TT left-to-right), with local ceilings
$n_i \leq \min(N_i,\, r_{i-1} r_i)$ (singletons),
$r_i \leq r_{i-1} n_i$, $r_{i-1} \leq n_i r_i$.

\begin{proposition}[Group ceiling]\label{prop:ceiling}
For a group $g$,
\begin{equation}\label{eq:ceiling}
n_g \;\leq\;
\min\Bigl( N_g,\ \sum_{i \in g} \min\bigl( N_g,\ r_{i-1}\, r_i \bigr) \Bigr).
\end{equation}
\end{proposition}

\begin{proof}
By Theorem~\ref{thm:sharedrank} and \eqref{eq:skeleton},
$\rank[\,\T_{(i_1)}|\cdots|\T_{(i_k)}\,]
\leq \min(N_g,\ \sum_i \rank \T_{(i)})
\leq \min(N_g,\ \sum_i \min(N_g,\ r_{i-1} r_i))$.
\end{proof}

A shared basis column is useless only if useless for \emph{every} mode of the
group: the per-mode ceilings \emph{add} across the group rather than
min-reduce, and $n_g$ may exceed $r_{i-1} r_i$ at individual modes without
the representation being reducible. Ceiling~\eqref{eq:ceiling} couples $n_g$
to TT ranks at distant chain positions, so the constraint incidence graph
acquires cycles; generic monotone constraint propagation on cyclic graphs
requires fixed-point iteration, in contrast to the tree-structured unshared
case (dynamic programming on a tree; one ordered pass). Nevertheless:

\begin{lemma}[Self-protection]\label{lem:selfprotect}
After the cap
$n_g \leftarrow \min\bigl(N_g,\ \sum_{i\in g} \min(N_g,\, r_{i-1}r_i)\bigr)$,
every TT constraint at the group's own modes involving $n_g$ holds
automatically, given that all $n_g$ values used in the TT phases satisfied
$n_g \leq N_g$ (phase~1).
\end{lemma}

\begin{proof}
The sum dominates each of its own terms, so
$n_g \geq \min(N_g,\, r_{i-1} r_i)$ for each $i \in g$.
\emph{Case $r_{i-1} r_i \leq N_g$:} then $n_g \geq r_{i-1} r_i$, and since
ranks are $\geq 1$,
\[
r_{i-1}\, n_g \;\geq\; r_{i-1}^2\, r_i \;\geq\; r_i,
\qquad
n_g\, r_i \;\geq\; r_{i-1}\, r_i^2 \;\geq\; r_{i-1},
\]
so both the left- and right-unfolding constraints hold.
\emph{Case $r_{i-1} r_i > N_g$:} the $i$-th term equals $N_g$, hence the sum
is $\geq N_g$ and $n_g = N_g$ exactly; the constraints
$r_i \leq r_{i-1} N_g$ and $r_{i-1} \leq N_g\, r_i$ are then precisely those
enforced by phase~1 ($n \leq N$) together with the standard TT phases, since
every $n_g$ value used there was $\leq N_g$.
\end{proof}

\begin{theorem}[Single pass]\label{thm:singlepass}
The three-phase sweep with (a) group-wide Tucker reductions, (b) the
ceiling~\eqref{eq:ceiling}, and (c) re-evaluation and re-capping of $n_g$ at
every group-mode visit during the left-to-right phase (standard
Tucker-then-TT ordering at each core), reaches the componentwise-greatest
nondegenerate rank vector below the proposal in a single pass.
\end{theorem}

\begin{proof}
Monotone cap iteration from the proposal converges to the greatest fixed
point (Knaster--Tarski/Kleene on the finite lattice of integer rank vectors);
it remains to show one ordered pass attains it, i.e.\ that no constraint is
violated at exit. Constraints not involving any $n_g$ are handled exactly as
in the unshared tree argument. For constraints involving $n_g$:

(i) \emph{Staleness within the sweep.} $n_g$ is monotone nonincreasing
through the sweep and $\leq N_g$ from phase~1 onward. A TT cap at group mode
$i$ was enforced with the then-current value; the left-to-right phase touches
$r_i$ for the last time at position $i$, before any later re-evaluation can
reduce $n_g$; final consistency at mode $i$ then follows from
Lemma~\ref{lem:selfprotect} applied with final values, since
$n_g^{\mathrm{final}} \geq \min(N_g,\, r_{i-1} r_i)$ with the final (frozen)
$r_{i-1}, r_i$.

(ii) \emph{Last group mode.} The ceiling at the group's last mode $i_k$ is
evaluated before that mode's own TT cap; if the cap binds
($r_{i_k} \leftarrow r_{i_k-1} n_g$), the mode's ceiling term becomes
$\min(N_g,\, r_{i_k-1}^2\, n_g) \geq \min(N_g,\, n_g) = n_g$, so the just-set
$n_g$ remains consistent.

(iii) \emph{No propagation.} $n_g$ appears in no constraint outside group
$g$'s modes, and by Lemma~\ref{lem:selfprotect} ceiling-driven reductions of
$n_g$ never force TT reductions. Hence the direction
$r \to n_g$ flows while $n_g \to r$ is vacuous; every cycle of the constraint
incidence graph contains a vacuous edge and is never traversed, and no
cross-group cascade exists.

At exit, all constraints hold, so the pass output is a fixed point; since
each cap is the loosest reduction consistent with its constraint and caps are
monotone, the output is the greatest fixed point below the proposal.
\end{proof}

\begin{remark}
The proof uses the specific phase structure: phase~1 ($n \leq N$)
\emph{precedes} the TT phases, and the left-to-right phase applies the Tucker
step before the TT step at each core. Implementations should assert the
theorem empirically (a second pass changes nothing) rather than assume it
survives reorderings. Note also that sharing only ever \emph{relaxes} Tucker
ceilings relative to per-mode removal, since the sum in \eqref{eq:ceiling}
dominates each per-mode term.
\end{remark}

\begin{thebibliography}{9}

\bibitem{MR2026}
A.~Molozhavenko and M.~Rakhuba,
\emph{Optimization on the extended tensor-train manifold with shared
factors}, Computational and Applied Mathematics \textbf{45}:221, 2026.

\bibitem{Pesh2024}
I.~Peshekhonov, A.~Arzhantsev, and M.~Rakhuba,
\emph{Training a Tucker model with shared factors: a Riemannian optimization
approach}, AISTATS, PMLR 238, 2024.

\bibitem{KSV2014}
D.~Kressner, M.~Steinlechner, and B.~Vandereycken,
\emph{Low-rank tensor completion by Riemannian optimization},
BIT Numerical Mathematics \textbf{54}, 447--468, 2014.

\bibitem{T4S}
N.~Alger, B.~Christierson, P.~Chen, and O.~Ghattas,
\emph{Tucker tensor train Taylor series}, arXiv:2603.21141, 2026.

\end{thebibliography}

\end{document}
